%=====================================================================
\chapter{Local Search and Optimisation}\label{ch:local}
%=====================================================================
\locator{2}{Part II --- Problem Solving by Search}

\begin{outcomes}
By the end of this chapter you will be able to:
\begin{tight}
  \item \textbf{Recognise} the problems for which local search is appropriate ---
        those where the goal is a good final state and the path to it is of no
        interest.
  \item \textbf{Implement} steepest-ascent hill climbing, and \textbf{diagnose}
        the three landscape features that defeat it: local maxima, plateaus and
        ridges.
  \item \textbf{Explain} simulated annealing's acceptance rule, and \textbf{say}
        what the temperature schedule controls at each end.
  \item \textbf{Describe} a genetic algorithm's four operators, and
        \textbf{identify} the assumption about the problem that crossover relies
        on.
  \item \textbf{Compare} local search methods at an \emph{equal evaluation budget},
        and \textbf{explain} why any other comparison is meaningless.
\end{tight}
\end{outcomes}

\begin{prereq}
\chref{uninformed} and \chref{informed} (the systematic searches this chapter
abandons, and the objective-function thinking it keeps). \chref{python} for
representing a state as a tuple. This chapter needs no explored set and no
frontier, and that is the point.
\end{prereq}

\begin{keyterms}
local search $\bullet$ objective function $\bullet$ state-space landscape $\bullet$
global maximum $\bullet$ local maximum $\bullet$ plateau $\bullet$ shoulder
$\bullet$ ridge $\bullet$ hill climbing $\bullet$ steepest ascent $\bullet$
sideways move $\bullet$ random-restart hill climbing $\bullet$ simulated annealing
$\bullet$ temperature schedule $\bullet$ local beam search $\bullet$ genetic
algorithm $\bullet$ selection $\bullet$ crossover $\bullet$ mutation $\bullet$
evaluation budget
\end{keyterms}

\section{When the Path Does Not Matter}

Every search so far returned a \emph{path}: do this, then this, then this. For the
release plan that was exactly right --- the order is the deliverable.

Now a different question. Nine developers must be assigned to four workstreams.
Each developer is better at some streams than others, and each stream wants a
particular headcount. There is no sequence here; there is only an assignment, and it
is either good or bad. Nobody cares in what order you decided it.

For problems of that shape, keeping a frontier and an explored set is pure waste.
\term{Local search} keeps a single current state and tries to improve it, using no
memory of how it got there. It gives up completeness and optimality, and in exchange
it handles spaces that systematic search cannot approach at all.

\begin{examplebox}[The staffing problem]
Nine developers, four streams (\texttt{api}, \texttt{ui}, \texttt{data},
\texttt{infra}). Each developer has a skill from $1$ to $5$ on each stream. The
target headcount is $(3,2,2,2)$.

A state is a tuple of nine stream indices, e.g.\ $(0,0,1,2,3,2,1,3,0)$. The
objective to maximise is
\[
  \mathrm{score}(s) \;=\;
  \underbrace{\sum_{i} \mathrm{skill}(d_{i}, s_{i})}_{\text{skill match}}
  \;-\;
  4 \sum_{j} \bigl| \mathrm{count}_{j}(s) - \mathrm{target}_{j} \bigr|
\]
The space has $4^{9} = 262{,}144$ states, small enough that the lab can find the
true optimum by exhaustive enumeration --- it is $\mathbf{42}$ --- and then check
how close each method gets. That is the only reason the problem is this small.
Twelve developers and six streams would be $6^{12} \approx 2.2$ billion.
\end{examplebox}

\section{The State-Space Landscape}

Picture the objective as a surface over the space of states: every state is a point
on the floor and its score is the height. Local search is a climber on that surface
who can see only the immediately adjacent ground. \dsref{landscape} names the
features that decide the outcome.

\dsfig{%
\begin{tikzpicture}[font=\scriptsize, x=1.25cm, y=1.0cm]

% --- axes ------------------------------------------------------------
\draw[gray!55, line width=0.8pt, -{Stealth[length=1.8mm]}] (0,0) -- (0,5.4)
  node[above, font=\scriptsize, text=neutral] {objective};
\draw[gray!55, line width=0.8pt, -{Stealth[length=1.8mm]}] (0,0) -- (9.9,0)
  node[right, font=\scriptsize, text=neutral, align=left] {state space};

% --- the landscape ---------------------------------------------------
\draw[primary, line width=1.3pt, smooth]
  plot coordinates {
  (0,1.2) (0.6,2.4) (1.2,2.9) (1.8,2.5) (2.4,1.6) (3.0,1.9)
  (3.6,2.6) (4.2,2.6) (4.8,2.6) (5.4,2.6) (6.0,3.4) (6.6,4.6)
  (7.2,4.9) (7.8,4.3) (8.4,2.8) (9.0,2.0) (9.6,2.3)};

% --- global maximum --------------------------------------------------
\draw[greenhl, line width=1pt, -{Stealth[length=1.8mm]}] (7.2,5.35) -- (7.2,5.05);
\node[font=\scriptsize\bfseries, text=greenhl!50!black, anchor=south]
  at (7.2,5.35) {global maximum};

% --- local maximum ---------------------------------------------------
\draw[accent, line width=1pt, -{Stealth[length=1.8mm]}] (1.2,3.95) -- (1.2,3.1);
\node[font=\scriptsize\bfseries, text=accent!55!black, anchor=south]
  at (1.2,3.95) {local maximum};
\node[font=\tiny\itshape, text=accent!55!black, anchor=north, align=center]
  at (1.2,1.05) {hill climbing};
\node[font=\tiny\itshape, text=accent!55!black, anchor=north, align=center]
  at (1.2,0.72) {stops here};

% --- shoulder --------------------------------------------------------
\draw[goldhl, line width=1pt, decorate,
      decoration={brace, amplitude=4pt, mirror}] (3.6,2.45) -- (5.4,2.45);
\node[font=\scriptsize\bfseries, text=goldhl!55!black, anchor=north]
  at (4.5,2.25) {shoulder: flat, but it leads up};

% --- the dip ---------------------------------------------------------
\node[font=\tiny\itshape, text=neutral, anchor=north, align=center]
  at (2.42,1.42) {the dip that};
\node[font=\tiny\itshape, text=neutral, anchor=north, align=center]
  at (2.42,1.09) {stops a climber};

\end{tikzpicture}}%
{The state-space landscape. Local search sees only the ground next to it. A
\emph{local maximum} has no better neighbour but is not the best state; a
\emph{plateau} is flat, so there is no gradient to follow, and a \emph{shoulder} is
a plateau that does lead upward --- which is why forbidding flat moves loses
solutions and allowing them without limit loops forever. A \emph{ridge} is a run of
local maxima along which no single move goes up, though a combination of two
would.}{landscape}

\section{Hill Climbing}

Look at every neighbour, move to the best one, stop when no neighbour is better.
That is \term{steepest-ascent hill climbing}, and it is about eight lines long.

It is also, on this problem, close to useless: from a random start it reaches the
optimum of $42$ in roughly one run in two hundred, and its average final score is
$30.3$. The landscape is full of local maxima because the headcount penalty and the
skill match pull against each other --- moving a developer to the stream she is best
at often unbalances two streams at once, and no \emph{single} move repairs it.

That is a ridge, and it is the characteristic failure. Hill climbing is not a bad
implementation of a good idea; it is a correct implementation of an idea that does
not work unaided.

\begin{alertbox}[The plateau dilemma]
If hill climbing stops when no neighbour is \emph{strictly} better, it halts on
every shoulder --- flat ground that would have led upward. If it allows sideways
moves onto equal-scoring neighbours, it can circle a genuine plateau forever.

The standard compromise is to allow sideways moves up to a fixed limit (a hundred,
say) and then stop. It is a hack, it works, and it is worth knowing that this is
what the literature does rather than assuming there is something more elegant.
\end{alertbox}

\section{Three Ways Out}

\subsection{Random-restart hill climbing}

Run hill climbing from a random start; keep the best result; repeat until the budget
is gone. It is three extra lines, it has no parameters, it is trivially parallel, and
it is a far stronger baseline than its reputation suggests --- on this problem it
lifts the mean from $30.3$ to $39.5$ and finds the optimum in $37$ runs out of $200$.

Remember that figure while reading the rest of the chapter. Any more sophisticated
method must beat random restart \emph{at the same budget}, or it is not earning its
complexity.

\subsection{Simulated annealing}

Instead of always climbing, pick a random neighbour and accept it if it is an
improvement --- but if it is worse, accept it anyway with probability
\[
  p \;=\; \exp\!\left(\frac{\Delta}{T}\right),
  \qquad \Delta = \mathrm{score}(\text{new}) - \mathrm{score}(\text{current}) < 0,
\]
where $T$ is a \term{temperature} that starts high and is lowered to zero on a
schedule. The name is from metallurgy: cool a metal slowly and its atoms settle into
a low-energy crystal; cool it fast and they freeze into whatever disordered state
they happened to be in.

Read the formula at its two extremes. When $T$ is large, $\Delta/T$ is near zero and
$p$ is near $1$ --- almost anything is accepted, and the search wanders freely,
ignoring the landscape. When $T$ approaches zero, $p$ approaches zero for any
negative $\Delta$, and the method becomes hill climbing. The schedule therefore
interpolates between a random walk and a greedy climb, and the art is all in how
fast. \dsref{annealing} plots the acceptance probability.

\dsfig{%
\begin{tikzpicture}
\begin{axis}[
  width=11.8cm, height=5.6cm,
  xlabel={how much worse the candidate is, $-\Delta$},
  ylabel={probability of accepting it},
  xlabel style={font=\scriptsize}, ylabel style={font=\scriptsize},
  tick label style={font=\tiny},
  xmin=0, xmax=10, ymin=0, ymax=1.05,
  grid=both, major grid style={gray!18}, minor grid style={gray!8},
  legend style={font=\tiny, at={(0.97,0.97)}, anchor=north east,
                draw=gray!40, fill=white},
  legend cell align=left, domain=0:10, samples=80]

\addplot[accent, line width=1.2pt]   {exp(-x/8)};
\addlegendentry{$T=8$ (start: almost anything goes)}
\addplot[goldhl, line width=1.2pt]   {exp(-x/4)};
\addlegendentry{$T=4$}
\addplot[secondary, line width=1.2pt]{exp(-x/2)};
\addlegendentry{$T=2$}
\addplot[primary, line width=1.2pt]  {exp(-x/0.5)};
\addlegendentry{$T=0.5$ (end: effectively hill climbing)}

\node[font=\tiny\itshape, text=neutral, anchor=west, align=left]
  at (axis cs:4.6,0.78)
  {a big step downhill is\\ accepted early, rarely late};
\end{axis}
\end{tikzpicture}}%
{Simulated annealing's acceptance probability $\exp(\Delta/T)$. High temperature
accepts almost any move; as $T$ falls the method converges on hill climbing. Note
that an \emph{improving} move is always accepted, at every
temperature.}{annealing}

There is a theoretical result worth knowing and not over-trusting: if the
temperature is lowered slowly enough, simulated annealing finds the global optimum
with probability approaching one. ``Slowly enough'' is slower than enumerating the
space, so the guarantee is of no operational use. What matters in practice is that
the method is robust, needs one parameter, and does well.

\subsection{Local beam search}

Keep $k$ states rather than one; generate all their successors; keep the best $k$.
This is not the same as $k$ parallel hill climbs, because the $k$ slots are
reallocated globally --- useful information passes between the threads, and the
population tends to concentrate in the promising region. With randomised selection
of which successors to keep, it is \term{stochastic beam search}, and it is the
direct ancestor of the next idea.

\section{Genetic Algorithms}

A genetic algorithm is stochastic beam search with one extra operator: states are
strings, and new states are made by \emph{combining two parents} rather than
mutating one. Four operators:

\begin{enumerate}[leftmargin=2.2em, itemsep=2pt, topsep=4pt]
  \item \textbf{Selection} --- choose parents with a bias toward high scores. The
        lab uses tournament selection: sample three at random, keep the best.
  \item \textbf{Crossover} --- cut both parents at the same random point and splice.
        With nine developers, cut after developer $4$ and take the first four
        assignments from one parent and the last five from the other.
  \item \textbf{Mutation} --- with small probability, reassign a developer at
        random. This is what stops the population collapsing to one state.
  \item \textbf{Elitism} --- carry the best few unchanged into the next generation,
        so a good solution is never lost.
\end{enumerate}

\begin{conceptbox}
Crossover is the whole claim of a genetic algorithm, and it rests on an assumption:
that the string decomposes into \emph{blocks that are independently good}. Splicing
the first four developers' assignments from one parent onto the last five from
another helps only if those are meaningful units that can be judged apart from each
other.

On the staffing problem the assumption holds reasonably well, and it is worth seeing
why. Most of a developer's contribution is her own skill score, which depends on her
assignment alone --- that part decomposes perfectly. The headcount penalty is the
only coupling, and it is a mild global correction rather than a pairwise tangle. So
inherited blocks mostly keep their value, and the measurement below shows the
genetic algorithm winning clearly.

That is a happy outcome and not a general law. On a problem where the objective
depends on \emph{pairs} of positions --- who is rostered alongside whom, which two
services must not share a host --- crossover routinely destroys exactly the
structure it inherits, and the same code performs badly. Crossover is a claim about
your problem. Check it; do not assume it.
\end{conceptbox}

\section{Comparing Them Honestly}

Almost every published comparison of these methods is worthless, for one reason: the
methods are given different amounts of work. A genetic algorithm with a population
of $100$ over $500$ generations has evaluated the objective $50{,}000$ times; hill
climbing from one start evaluates it a few hundred times. Declaring the former
better is not a finding.

The only fair comparison fixes the \term{evaluation budget}: every method may call
the objective function the same number of times. \dsref{local-compare} and
Table~\ref{tab:local} do that at $4{,}000$ calls, over $200$ independent runs with
seeds $0$ to $199$.

\dsfig{%
\begin{tikzpicture}
\begin{axis}[
  width=12.6cm, height=5.2cm,
  xbar, y=0.72cm,
  xmin=0, xmax=46,
  xlabel={mean final score over 200 runs (the optimum is 42)},
  xlabel style={font=\scriptsize},
  tick label style={font=\scriptsize},
  ytick=data,
  yticklabels={Hill climbing (one start), Random restart,
               Simulated annealing, Genetic algorithm},
  yticklabel style={align=right, font=\scriptsize},
  nodes near coords, nodes near coords style={font=\tiny, text=neutral},
  point meta=explicit symbolic,
  enlarge y limits=0.18,
  axis lines=left, xmajorgrids, grid style={gray!18}]

\addplot[draw=primary!70, fill=primary!20] coordinates {
  (30.34,0) [30.3]
  (39.45,1) [39.5]
  (39.99,2) [40.0]
  (41.54,3) [41.5]};

\draw[greenhl, line width=1.1pt, dash pattern=on 3pt off 2pt]
  (axis cs:42,-0.6) -- (axis cs:42,3.6);
\node[font=\tiny\bfseries, text=greenhl!50!black, anchor=south, rotate=90]
  at (axis cs:42.7,0.6) {optimum 42};
\end{axis}
\end{tikzpicture}}%
{Mean final score at an equal budget of $4{,}000$ objective evaluations per run.
Hill climbing alone is far behind; the three escape strategies are close together,
with the genetic algorithm ahead.}{local-compare}

\smallskip
\noindent\begin{longtable}{@{}L{4.3cm}L{2.5cm}L{2.2cm}L{2.0cm}L{2.4cm}@{}}
\toprule
\textbf{Method} & \textbf{Found optimum} & \textbf{Mean score} &
\textbf{Worst run} & \textbf{Parameters} \\
\midrule
\endfirsthead
\toprule
\textbf{Method} & \textbf{Found optimum} & \textbf{Mean score} &
\textbf{Worst run} & \textbf{Parameters} \\
\midrule
\endhead
Hill climbing, one start & 1 / 200 & 30.34 & 22 & none \\
Random-restart hill climbing & 37 / 200 & 39.45 & 36 & none \\
Simulated annealing & 51 / 200 & 39.99 & 37 & initial temperature \\
Genetic algorithm & \textbf{160 / 200} & \textbf{41.54} & 37 &
population, mutation rate, tournament size \\
\bottomrule
\caption{Local search on the staffing problem at $4{,}000$ evaluations per run,
$200$ runs, seeds $0$--$199$. The exhaustive optimum is $42$.}\label{tab:local}
\end{longtable}

Four things to take from the table.

\textbf{Escaping local maxima is what matters, and almost any escape works.} The gap
from hill climbing to the \emph{weakest} escape strategy is $9.1$ points of mean
score; the gap from the weakest escape to the best is $2.1$. Nearly all the available
improvement comes from not being a pure hill climber, and very little from which
alternative you pick.

\textbf{Random restart is a serious baseline.} No parameters, three lines, perfect
parallelism, and it comes within $2.1$ points of the winner. Measure it before
reaching for anything more interesting; if your sophisticated method cannot beat it,
you have learned something useful cheaply.

\textbf{The genetic algorithm wins here, on a problem whose structure suits it.} It
reached the optimum in four runs out of five. Section~\ref{sec:accounting} shows how
easily that result can be reversed by an apparently innocent change to how the
budget is counted --- which is the real lesson of this section.

\textbf{None of them is reliable on a single run.} Even the winner missed the optimum
in $40$ of $200$ runs. If you need \emph{the} optimum rather than a good answer,
local search is the wrong tool and \chref{csp} is the right one.

\subsection{How you count an evaluation can reverse the ranking}
\label{sec:accounting}

\begin{alertbox}[The same four methods, two defensible accountings]
A genetic algorithm evaluates fitness in two places: once per individual per
generation, and again inside each tournament when it compares candidate parents. The
standard --- and correct --- practice is to evaluate each individual \emph{once} per
generation and cache the value for every selection decision, so a generation costs
exactly \texttt{pop} evaluations. That is what Table~\ref{tab:local} does, and the
genetic algorithm comes first with a mean of $41.54$.

Charge each tournament comparison as a fresh evaluation instead --- two tournaments
of three per child, so roughly six times as many --- and the same code, same seeds,
same budget, produces a mean of $38.25$ and $29/200$ optima. \emph{The genetic
algorithm moves from first place to last.} Nothing about the algorithm changed; only
the bookkeeping did.

\smallskip
Which accounting is right? Caching is right, because in any real application the
objective is the expensive part and no competent implementation would recompute a
fitness it already holds. But notice how little the question announces itself, and
how completely the answer changes the published conclusion. When you read a
comparison of metaheuristics, the first question is not which method won --- it is
what the authors counted.
\end{alertbox}

\begin{worked}[Why one developer cannot be moved to fix a ridge]
Take the state in which every developer is assigned to the stream she is
individually best at. Reading the skill table: Ayesha $\to$ \texttt{api} ($5$),
Bilal $\to$ \texttt{api} ($4$), Chen $\to$ \texttt{ui} ($5$), Dina $\to$
\texttt{data} ($5$), Emre $\to$ \texttt{infra} ($5$), Farah $\to$ \texttt{data}
($5$), Gul $\to$ \texttt{ui} ($5$), Hassan $\to$ \texttt{infra} ($5$), Iqra $\to$
tied at $3$, say \texttt{api}.

\smallskip
\textbf{Skill match:} $5+4+5+5+5+5+5+5+3 = 42$.
\textbf{Headcount:} \texttt{api} $3$, \texttt{ui} $2$, \texttt{data} $2$,
\texttt{infra} $2$ --- exactly the target $(3,2,2,2)$, so the imbalance penalty is
zero.
\textbf{Score:} $42 - 0 = 42$. This is the global optimum, and the lab confirms it
by enumeration.

\smallskip
\textbf{Now start one move away from it} --- swap Iqra to \texttt{ui}. Headcount
becomes $(2,3,2,2)$: two streams are off target by one each, so the imbalance is
$2$ and the penalty is $8$. Skill match becomes $42-3+3 = 42$. Score $34$.

\smallskip
\textbf{The trap.} From that state, consider every single move. Moving Iqra back to
\texttt{api} recovers $42$ --- so this particular state is \emph{not} a local
maximum, and hill climbing escapes it. The genuine traps are states where two
developers are each in the wrong stream in a way that balances the headcount. Move
either one alone and the headcount goes off target, costing $8$, while the skill
gain is at most $4$. So both single moves go \emph{down}, and only the simultaneous
pair goes up.

\smallskip
\textbf{That is a ridge}, and it is why hill climbing reaches $42$ only once in two
hundred starts. It also explains why simulated annealing helps: accepting the first
of the two moves, at a cost of $4$, is exactly the kind of temporary worsening that
a positive temperature permits and a hill climber forbids.
\end{worked}

\begin{pitfall}
\begin{tight}
  \item \textbf{Comparing methods at unequal budgets.} The most common error in this
        area, and it always flatters the method with the largest population. Fix the
        number of objective evaluations.
  \item \textbf{Using local search when you need a guarantee.} It gives a good
        state, never a proof. ``No engineer is double-booked'' is not something to
        establish by annealing --- use \chref{csp}.
  \item \textbf{Allowing unlimited sideways moves.} An infinite loop on a plateau.
        Cap them.
  \item \textbf{Reaching for a genetic algorithm first.} It has the most parameters
        and the strongest assumption --- that crossover recombines meaningful blocks.
        It won here, but only $2.1$ mean points ahead of parameterless random
        restart, and on a problem with pairwise structure it would lose. Establish
        the baseline first.
  \item \textbf{Not stating what an ``evaluation'' means in a comparison.} Charging
        a genetic algorithm per tournament comparison instead of per individual per
        generation moved it from first place to last on this problem, with no change
        to the code (Section~\ref{sec:accounting}).
  \item \textbf{An annealing schedule that cools too fast.} At $T \approx 0$ from
        the second iteration it \emph{is} hill climbing, with extra code. Print the
        acceptance rate over time; if it is near zero early, the schedule is wrong.
  \item \textbf{Reporting the best of many runs as the method's performance.} Report
        the mean and the worst as well, or you are describing your luck.
\end{tight}
\end{pitfall}

%---------------------------------------------------------------------
\begin{lab}[Staffing by Annealing and by Evolution]
Four methods, one objective, one budget. The lab finds the true optimum by
enumeration first, so every method can be scored against the truth rather than
against each other.
\begin{lstlisting}[language=Python]
"""Chapter 6 lab -- local search on the staffing problem.

Nine developers, four workstreams. Reproduces Table 6.1. Every method
gets exactly the same number of objective-function evaluations, which
is the only comparison that means anything.
"""

import itertools
import math
import random

DEVS = ("Ayesha", "Bilal", "Chen", "Dina", "Emre",
        "Farah", "Gul", "Hassan", "Iqra")
STREAMS = ("api", "ui", "data", "infra")

# skill of each developer on each stream, 1 (weak) to 5 (strong)
SKILL = {
    "Ayesha": (5, 2, 3, 1),
    "Bilal":  (4, 3, 2, 2),
    "Chen":   (2, 5, 1, 2),
    "Dina":   (1, 4, 5, 2),
    "Emre":   (3, 1, 4, 5),
    "Farah":  (2, 2, 5, 3),
    "Gul":    (1, 5, 2, 3),
    "Hassan": (4, 1, 2, 5),
    "Iqra":   (3, 3, 3, 3),
}
TARGET = (3, 2, 2, 2)      # headcount wanted per stream
PENALTY = 4                # cost per developer away from target
NS = len(STREAMS)
BUDGET = 4000              # objective evaluations allowed per run


def raw_score(state):
    fit = sum(SKILL[DEVS[i]][s] for i, s in enumerate(state))
    counts = [0] * NS
    for s in state:
        counts[s] += 1
    off = sum(abs(counts[j] - TARGET[j]) for j in range(NS))
    return fit - PENALTY * off


class Budget:
    """Counts objective evaluations, so no method can cheat."""

    def __init__(self, cap):
        self.cap, self.n = cap, 0

    def score(self, state):
        self.n += 1
        return raw_score(state)

    def left(self):
        return self.n < self.cap


def neighbours(state):
    for i in range(len(state)):
        for s in range(NS):
            if s != state[i]:
                yield state[:i] + (s,) + state[i + 1:]


def random_state(rng):
    return tuple(rng.randrange(NS) for _ in DEVS)


def one_neighbour(state, rng):
    i = rng.randrange(len(state))
    s = rng.randrange(NS)
    while s == state[i]:
        s = rng.randrange(NS)
    return state[:i] + (s,) + state[i + 1:]


# --- 1. steepest-ascent hill climbing ------------------------------
def hill_climb(b, rng, start=None):
    st = start if start is not None else random_state(rng)
    while b.left():
        cur = b.score(st)
        best, bestv = None, cur
        for n in neighbours(st):
            if not b.left():
                break
            v = b.score(n)
            if v > bestv:
                best, bestv = n, v
        if best is None:
            return st, cur          # no better neighbour: stuck
        st = best
    return st, raw_score(st)


# --- 2. random-restart hill climbing -------------------------------
def random_restart(b, rng):
    best, bestv = None, None
    while b.left():
        st, v = hill_climb(b, rng)
        if bestv is None or v > bestv:
            best, bestv = st, v
    return best, bestv


# --- 3. simulated annealing ----------------------------------------
def anneal(b, rng, t0=8.0):
    st = random_state(rng)
    cur = b.score(st)
    best, bestv = st, cur
    while b.left():
        T = t0 * (1 - b.n / b.cap) + 1e-9      # linear cooling
        cand = one_neighbour(st, rng)
        v = b.score(cand)
        d = v - cur
        if d > 0 or rng.random() < math.exp(d / T):
            st, cur = cand, v
            if cur > bestv:
                best, bestv = st, cur
    return best, bestv


# --- 4. genetic algorithm ------------------------------------------
def genetic(b, rng, pop=30, mut=0.12):
    """Standard fitness accounting: every individual is evaluated once
    per generation and the value is then cached for every selection
    decision in that generation. So one generation costs exactly `pop`
    evaluations -- no selection is ever free.

    This matters more than it looks. Charging each tournament
    comparison separately instead costs about six times as much per
    generation, and the method drops from first place to last. See the
    alertbox in Section 6.6.
    """
    P = [random_state(rng) for _ in range(pop)]
    bestx, bestv = None, None
    while b.n + pop <= b.cap:
        fit = [b.score(x) for x in P]          # pop evaluations, once
        order = sorted(range(pop), key=lambda i: -fit[i])
        if bestv is None or fit[order[0]] > bestv:
            bestv, bestx = fit[order[0]], P[order[0]]
        elite = [P[i] for i in order[: pop // 5]]
        nxt = list(elite)

        def tournament():
            """Best of three, judged on the cached fitness."""
            picked = rng.sample(range(pop), 3)
            return P[max(picked, key=lambda i: fit[i])]

        while len(nxt) < pop:
            a, c = tournament(), tournament()
            cut = rng.randrange(1, len(DEVS))
            child = a[:cut] + c[cut:]
            child = tuple(rng.randrange(NS) if rng.random() < mut else g
                          for g in child)
            nxt.append(child)
        P = nxt
    return bestx, bestv


def exhaustive():
    """4^9 = 262144 states: small enough to know the truth."""
    best, arg = None, None
    for st in itertools.product(range(NS), repeat=len(DEVS)):
        v = raw_score(st)
        if best is None or v > best:
            best, arg = v, st
    return best, arg


def show(state):
    by = {j: [] for j in range(NS)}
    for i, s in enumerate(state):
        by[s].append(DEVS[i])
    return "  ".join("%s: %s" % (STREAMS[j], ", ".join(by[j]))
                     for j in range(NS))


if __name__ == "__main__":
    opt, arg = exhaustive()
    print("states: %d      exhaustive optimum: %d" % (NS ** len(DEVS), opt))
    print(show(arg))
    print()
    print("budget: %d evaluations per run, 200 runs, seeds 0-199" % BUDGET)
    print()
    print("method                        optimum    mean   worst")
    print("-" * 56)

    methods = (("hill climbing, one start", hill_climb),
               ("random restart          ", random_restart),
               ("simulated annealing     ", anneal),
               ("genetic algorithm       ", genetic))
    table = {}
    for name, fn in methods:
        hits, total, worst = 0, 0, None
        for seed in range(200):
            rng = random.Random(seed)
            _, v = fn(Budget(BUDGET), rng)
            total += v
            worst = v if worst is None else min(worst, v)
            if v == opt:
                hits += 1
        table[name.strip()] = (hits, total / 200, worst)
        print("%s  %4d/200 %7.2f %6d"
              % (name, hits, total / 200, worst))

    hc = table["hill climbing, one start"]
    rr = table["random restart"]
    sa = table["simulated annealing"]
    ga = table["genetic algorithm"]
    escapes = (rr, sa, ga)

    # 1. plain hill climbing is hopeless here: the landscape is a ridge
    assert hc[0] <= 5
    # 2. every escape strategy is far better than no escape
    for m in escapes:
        assert m[1] > hc[1] + 5
    # 3. the gap to no-escape dwarfs the gap between the escapes
    spread = max(m[1] for m in escapes) - min(m[1] for m in escapes)
    gain = min(m[1] for m in escapes) - hc[1]
    assert spread < gain
    # 4. even the winner is unreliable on any single run
    assert max(m[0] for m in escapes) < 200

    print()
    print("escaping local maxima is worth %.1f points of mean score;"
          % gain)
    print("choosing the best escape is worth only %.1f more." % spread)
    print()
    print("random restart has no parameters at all and reaches %.2f."
          % rr[1])
    print("measure it before reaching for anything more interesting.")
    print()
    best = max(escapes, key=lambda m: m[1])
    print("the best method still missed the optimum in %d of 200 runs."
          % (200 - best[0]))
\end{lstlisting}
Two experiments to run afterwards. Set the annealing temperature \texttt{t0} to
$0.01$ and confirm that it degenerates into hill climbing --- the mean drops back
towards $30$. Then raise \texttt{PENALTY} from $4$ to $20$ and watch every method
get worse together: a harder penalty makes the ridges steeper, and no amount of
cleverness in the search compensates for a landscape that has been made hostile by
the modelling.
\end{lab}

%---------------------------------------------------------------------
\begin{checkpoint}
\begin{tightnum}
  \item Explain the difference between a plateau and a shoulder, and say what goes
        wrong if hill climbing (a)~forbids sideways moves and (b)~allows unlimited
        sideways moves.
  \item Simulated annealing is offered a neighbour that is $6$ points worse. Compute
        the acceptance probability at $T = 8$ and at $T = 0.5$, and say what the
        two numbers mean for the method's behaviour early and late in a run.
  \item Name the assumption crossover makes about a problem, say why the staffing
        objective largely satisfies it, and give one change to the objective that
        would break it.
\end{tightnum}
\end{checkpoint}

%---------------------------------------------------------------------
\begin{chaptersummary}
\begin{tight}
  \item \term{Local search} keeps one current state and no path. Use it when the
        answer is a good final state and the route to it is of no interest --- and
        never when you need a guarantee.
  \item The \term{state-space landscape} explains every failure: \term{local
        maxima}, \term{plateaus} (flat, possibly a \term{shoulder} that leads up)
        and \term{ridges} (no single move improves, though a combination would).
  \item \term{Hill climbing} moves to the best neighbour and stops when there is
        none. On the staffing problem it reaches the optimum once in $200$ starts,
        because the skill match and the headcount penalty form a ridge.
  \item \term{Random restart} costs three lines, has no parameters, parallelises
        perfectly and lifted the mean from $30.3$ to $39.5$. It is the baseline
        every other method must beat.
  \item \term{Simulated annealing} accepts a worsening move with probability
        $\exp(\Delta/T)$, interpolating from a random walk at high $T$ to hill
        climbing at $T \to 0$. One parameter, mean $40.0$, $51/200$ optima.
  \item A \term{genetic algorithm} adds \term{crossover} to stochastic beam search.
        Crossover assumes the state decomposes into independently good blocks; the
        staffing problem largely satisfies that, and the method led the table with
        $160/200$ optima and a mean of $41.5$.
  \item Compare methods only at an \emph{equal evaluation budget}, say exactly what
        an evaluation is, and report the mean and the worst run rather than the best.
        Charging the genetic algorithm per tournament comparison rather than per
        individual per generation moves it from first place to last on this problem
        without touching its code.
\end{tight}
\end{chaptersummary}

\begin{reviewq}
  \item \textbf{(MCQ)} Local search is appropriate when: (a)~the path to the goal is
        the answer (b)~only the final state matters (c)~optimality must be proved
        (d)~the space is small enough to enumerate.
  \item \textbf{(MCQ)} A \emph{ridge} defeats hill climbing because:
        (a)~all neighbours are equal (b)~no single move improves, though a
        combination would (c)~the objective is undefined (d)~the space is
        continuous.
  \item \textbf{(MCQ)} In simulated annealing, as $T \to 0$ the method becomes:
        (a)~a random walk (b)~breadth-first search (c)~hill climbing (d)~a genetic
        algorithm.
  \item \textbf{(MCQ)} The assumption behind crossover is that:
        (a)~the objective is linear (b)~the state decomposes into independently good
        blocks (c)~mutation is unnecessary (d)~the population is large.
  \item \textbf{(Short)} Why is comparing a genetic algorithm to hill climbing
        without fixing the evaluation budget meaningless? Describe the comparison
        you would run instead.
  \item \textbf{(Short)} Give the acceptance rule of simulated annealing and explain
        what each of its two extremes reduces to.
  \item \textbf{(Short)} The genetic algorithm led Table~\ref{tab:local} on the
        standard accounting and came last when each tournament comparison was charged
        separately. Say which accounting is correct and why, and what the episode
        implies about published comparisons of metaheuristics.
  \item \textbf{(Applied)} Iqra has skill $3$ on all four streams. Show that the
        optimal assignment puts her on \texttt{api}, by computing the score of
        moving her to each of the other three streams from the optimum. Then explain
        what her uniform skill row contributes to the difficulty of the landscape.
\end{reviewq}
